\section{总量一致的2026层级预测与经营动作} % src:交接/计划.json:论文结构[文件名=论文/6.层级预测与经营策略.tex].章节标题

\subsection{总量硬锚与层级份额收缩}

本问不再改动上一章已经通过封存测试的总量模型，而是把每个月的总销售额分配到区域、产品、客户和发货层级。对任一维度，记历史月份 $t$ 中层级 $g$ 的销售额为 $x_{g,t}$，总销售额为 $Y_t$，则历史份额 $s_{g,t}=x_{g,t}/Y_t$。预测阶段把长期份额 $L_g$、近窗份额 $R_g$ 与同月季节份额 $M_{g,m(t)}$ 收缩为
\begin{equation}
\begin{aligned}
    B_{g,t}&=\eta R_g+(1-\eta)M_{g,m(t)},\\
    \rho_g&=\frac{12}{12+\kappa(1+CV_g^2)},\\
    q_{g,t}&=\rho_gB_{g,t}+(1-\rho_g)L_g,\\
    \widehat s_{g,t}&=\frac{\max(q_{g,t},0)}{\sum_h\max(q_{h,t},0)}.
\end{aligned}
\end{equation} % src:求解/问题4/结果/2026层级预测_区域.json:份额收缩参数,份额分解诊断
其中，$\eta$ 控制近期信息相对同月季节信息的权重，$\kappa$ 控制向长期结构收缩的强度，$CV_g$ 是历史月度份额变异系数，$\rho_g$ 是局部信息保留系数。因为本题层级越细，份额波动和零销售月份越容易放大，$CV_g$ 增大时 $\rho_g$ 自动减小，偶然的大额订单便不会直接变成下一年的结构判断。

起初的做法是分别拟合区域、产品大类、产品子类、客户细分和发货模式的月度总额模型，后来发现各层级最多只有48个月，产品子类份额变异系数中位数为0.684、最大为2.147，单个子类最多出现16个零销售月。% src:交接/实验记录.json:[问题=问题4/层级预测方法选择].现象
这些现象意味着独立时序既会放大稀疏层级的方差，也无法自然回到同一个总量，因此改用总量硬锚加份额收缩。锁定参数取 $\eta=0.75$、$\kappa=12$，同月季节份额权重相应为0.25；参数只由2024年校准误差确定，未用2025年重新挑选。% src:求解/问题4/结果/份额收缩回测与灵敏度.json:锁定参数

层级点预测与可加总风险边界统一写为
\begin{equation}
    \widehat x_{g,t}=\widehat s_{g,t}\widehat Y_t,
    \qquad
    \widehat x^{(b)}_{g,t}=\widehat s_{g,t}\widehat Y_t^{(b)},
    \qquad b\in\{80L,80U,95L,95U\}.
\end{equation} % src:求解/问题4/结果/2026总量与一致性.json:指标含义,区间解释
其中，$\widehat Y_t$ 是问题三给出的月度总量点预测，$\widehat Y_t^{(b)}$ 是相应边界。归一化保证同一维度的份额非负且和为一，因而点预测和各条边界都严格加回总量；这种协调口径与层级预测强调的跨层加总约束一致\cite{hyndman2011hierarchical}。这些边界没有联合计入份额随机性，只表示给定预测份额后的可加总风险边界，不能解释为层级对象的完整概率区间。

产品层级还需要消除大类表与子类表之间的矛盾。令 $\widehat s^C_{k,t}$ 为大类 $k$ 的预测份额，$\widehat u_{j\mid k,t}$ 为子类 $j$ 在该大类内部的条件份额，则
\begin{equation}
    \widehat s^{SC}_{j,t}=\widehat s^C_{k,t}\widehat u_{j\mid k,t},
    \qquad
    \sum_{j\in k}\widehat s^{SC}_{j,t}=\widehat s^C_{k,t}.
\end{equation} % src:求解/问题4/结果/2026层级预测_产品层级.json:产品子类.建模口径,一致性检查.父子协调
其中，$\widehat u_{j\mid k,t}$ 在父级内部归一化。这样子类先加到所属大类，大类再加到月度总量，避免两套独立份额表给出互相冲突的经营口径。

\subsection{回测改进与2026年总量节奏} % src:求解/问题4/结果/2026总量与一致性.json:2026全年总量预测

固定复制去年同月份额虽然简单，却只保留单月结构。锁定后的收缩模型在2025年五维平均份额WAPE为29.015\%，固定份额基线为39.585\%，改善10.570个百分点；其中产品子类仍有44.803\%的份额WAPE，说明协调后的细层结果可用于方向与容量安排，不宜把小子类名次写成刚性承诺。% src:求解/问题4/结果/份额收缩回测与灵敏度.json:2025五维平均对比,2025伪实时回测.产品子类
近期权重与收缩强度的9组扰动使2025年平均份额WAPE落在29.010\%至31.031\%，全部低于同一基线。% src:求解/问题4/结果/份额收缩回测与灵敏度.json:关键参数灵敏度_2025仅作检验不改锁定参数,2025五维平均对比.固定去年同月份额基线平均份额WAPE_百分比

全年点预测为791951.82，95\%可加总风险边界为$[454803.67,1379541.60]$；季度销售额逐季抬升，第四季度占全年点预测的38.30\%，经营资源不能按四季平均摊派。% src:求解/问题4/结果/2026总量与一致性.json:2026全年总量预测,2026季度总量预测[季度=2026Q4]
问题三封存测试中，月度95\%风险带的实际覆盖率为83.33\%，低于名义水平；季度和全年边界又是月度上下边界分别相加的结果，没有单独校准联合覆盖率，故表\ref{tab:q4_quarter_total}只把它们用于资源压力测试，不作概率承诺。% src:求解/问题3/结果/2025评估指标.json:区间测试表现.95%覆盖率 % src:求解/问题3/结果/2026总量预测.json:接口说明.季度与全年区间

\begingroup
\setlength{\tabcolsep}{5pt}
\begin{longtable}{>{\centering\arraybackslash}p{0.16\textwidth}>{\centering\arraybackslash}p{0.22\textwidth}>{\centering\arraybackslash}p{0.24\textwidth}>{\centering\arraybackslash}p{0.26\textwidth}}
    \caption{2026年季度总量预测及可加总风险边界（销售额单位：万）}\label{tab:q4_quarter_total}\\ % src:求解/问题4/结果/2026总量与一致性.json:2026季度总量预测
    \toprule
    季度 & 点预测 & 80\%风险边界 & 95\%可加总风险边界\\
    \midrule
    2026Q1 & 12.35 & $[9.21,16.55]$ & $[7.00,21.77]$\\ % src:求解/问题4/结果/2026总量与一致性.json:2026季度总量预测[季度=2026Q1]
    2026Q2 & 15.82 & $[12.40,20.19]$ & $[9.44,26.51]$\\ % src:求解/问题4/结果/2026总量与一致性.json:2026季度总量预测[季度=2026Q2]
    2026Q3 & 20.70 & $[16.18,26.47]$ & $[11.78,36.37]$\\ % src:求解/问题4/结果/2026总量与一致性.json:2026季度总量预测[季度=2026Q3]
    2026Q4 & 30.33 & $[23.71,38.80]$ & $[17.26,53.31]$\\ % src:求解/问题4/结果/2026总量与一致性.json:2026季度总量预测[季度=2026Q4]
    \bottomrule
\end{longtable}
\endgroup

表\ref{tab:q4_quarter_total}中的旺季抬升在图\ref{fig:q4_fan_matrix}左侧表现为年末扇形带同步变宽。五个维度逐月汇总的最大误差为$5.82\times10^{-11}$，产品大类与子类的最大父子协调误差为$1.46\times10^{-11}$，均表明数值误差没有破坏经营口径。% src:求解/问题4/结果/2026总量与一致性.json:层级一致性检查

\begin{figure}[H]
    \centering
    \includegraphics[width=0.8\textwidth]{../求解/问题4/图片/2026预测扇形与策略机会—把握度矩阵.png}
    \caption{2026年总量风险带与层级机会—把握度矩阵} % src:交接/图注素材_B.json:[图名=2026预测扇形与策略机会—把握度矩阵]
    \label{fig:q4_fan_matrix}
\end{figure}

\subsection{机会—把握度评分与情景边界}

预测规模不能单独决定投入强度。对每个真实层级对象，机会分 $O_g$ 综合预测增量、预测规模和历史增长，把握度分 $C_g$ 综合问题二证据、份额稳定性和相对不确定性的逆序分位秩：
\begin{equation}
\begin{aligned}
    O_g&=0.45P_{\Delta,g}+0.35P_{S,g}+0.20P_{G,g},\\
    C_g&=0.45E_g+0.30P_{V,g}+0.25P_{U,g}.
\end{aligned}
\end{equation} % src:求解/问题4/结果/机会_把握度策略矩阵.json:评分公式
其中，$P_{\Delta,g},P_{S,g},P_{G,g}$ 分别是维度内增量、规模和历史增长的分位秩，$E_g$ 是双证据得分，$P_{V,g}$ 与 $P_{U,g}$ 分别奖励稳定份额和较低不确定性。机会分与把握度分均以60分划分高低，因而高销售容量但缺少驱动证据的对象仍可能落入小步试点区。% src:求解/问题4/结果/机会_把握度策略矩阵.json:评分公式.阈值

我们最初考虑按权重比例缩放已有总分，但发现总分无法还原机会分和把握度分的各三个组件，这种缩放会把口径偏移误当成对象翻转。% src:交接/实验记录.json:[问题=问题4/机会评分权重敏感性],[问题=问题4/把握度内部权重敏感性扩展]
因此从31个对象的原始分项重建分位秩，基准机会分与把握度分的最大复现误差均为0；固定机会阈值后，将四组机会权重、四组把握度权重与三档把握度阈值组合成48组全网格。% src:求解/问题4/结果/机会评分敏感性.json:输入说明,解释边界

全网格入选频率把稳定对象与临界对象直接分开：Technology、Accessories、Chairs、Phones、Storage、Tables和East在48组中全部入选；West和Binders分别入选44组和40组；Central仅入选21组，Consumer与Second Class均仅入选16组。% src:求解/问题4/结果/机会评分敏感性.json:全网格入选频率,跨网格汇总.各对象入选频率
全网格中核心清单规模介于8与17个对象之间，最低仅保留基准清单的57.14\%，故全入选对象可用于基准资源配置，未全入选对象必须等实际增量或容量信号触发后再追加。% src:求解/问题4/结果/机会评分敏感性.json:跨网格汇总

分类稳定不等于总量情景下的名次稳定。图\ref{fig:q4_fan_matrix}右侧的31个层级对象中，点预测与95\%下侧风险边界机会分的Spearman相关为0.701，前10项仅重合5项；West由点预测的第14位升至下侧风险边界的第1位。% src:求解/问题4/结果/机会_把握度策略矩阵.json:点预测与下界排序稳健性,策略矩阵[维度=区域,层级=West]
这一反转使West在悲观情景中承担基本盘防守，基准情景的新增资源则先投向East及全网格稳定产品。

\subsection{分层经营动作}

表\ref{tab:q4_actions}从稳定核心、高频入选和低频敏感对象中抽取代表，将基准销售额与悲观边界配对后再设触发动作。证据等级A表示全网格入选，B表示高频入选但存在翻转，C表示对权重或阈值较敏感；括号内同时给出真实入选次数，不将等级解释为概率。% src:求解/问题4/结果/机会评分敏感性.json:全网格入选频率

\begingroup
\setlength{\tabcolsep}{5pt}
\begin{longtable}{>{\centering\arraybackslash}p{0.13\textwidth}>{\centering\arraybackslash}p{0.14\textwidth}>{\centering\arraybackslash}p{0.14\textwidth}>{\centering\arraybackslash}p{0.28\textwidth}>{\centering\arraybackslash}p{0.15\textwidth}}
    \caption{代表对象的2026年销售额边界与触发动作（单位：万）}\label{tab:q4_actions}\\
    \toprule
    对象 & 基准销售额 & 悲观边界 & 触发动作 & 证据等级\\
    \midrule
    Phones & 11.50 & 6.60 & 超基准再扩容 & A（48/48）\\ % src:求解/问题4/结果/机会_把握度策略矩阵.json:策略矩阵[维度=产品子类,层级=Phones].2026点预测,2026_95%下界 % src:求解/问题4/结果/机会评分敏感性.json:全网格入选频率[对象标识=产品子类::Phones]
    Storage & 7.63 & 4.38 & 近下界收紧补货 & A（48/48）\\ % src:求解/问题4/结果/机会_把握度策略矩阵.json:策略矩阵[维度=产品子类,层级=Storage].2026点预测,2026_95%下界 % src:求解/问题4/结果/机会评分敏感性.json:全网格入选频率[对象标识=产品子类::Storage]
    East & 23.29 & 13.37 & 下行时仅留可撤预算 & A（48/48）\\ % src:求解/问题4/结果/机会_把握度策略矩阵.json:策略矩阵[维度=区域,层级=East].2026点预测,2026_95%下界 % src:求解/问题4/结果/机会评分敏感性.json:全网格入选频率[对象标识=区域::East]
    West & 25.69 & 14.75 & 悲观保供，超基准扩展 & B（44/48）\\ % src:求解/问题4/结果/机会_把握度策略矩阵.json:策略矩阵[维度=区域,层级=West].2026点预测,2026_95%下界 % src:求解/问题4/结果/机会评分敏感性.json:全网格入选频率[对象标识=区域::West]
    Central & 16.66 & 9.57 & 增量兑现后追加 & C（21/48）\\ % src:求解/问题4/结果/机会_把握度策略矩阵.json:策略矩阵[维度=区域,层级=Central].2026点预测,2026_95%下界 % src:求解/问题4/结果/机会评分敏感性.json:全网格入选频率[对象标识=区域::Central]
    \bottomrule
\end{longtable}
\endgroup

Phones的旺季上侧风险边界较点预测高约3.32万，这段容量区间适合检验仓储、采购额度和供应响应能否承接销售额波动，但数据没有销量与库存字段，不能换算成安全库存件数。% src:求解/问题4/结果/经营动作建议.json:动作清单[动作类型=库存管理,对象=Phones]
物流侧还核对了42条异常日期：年份修正与异常置缺口径的订单簇置换$p$值分别为0.108和0.052，两种处理都没有形成稳定驱动证据。% src:求解/问题4/结果/机会_把握度策略矩阵.json:物流双口径策略灵敏度
经营执行可据此设置三档口径：基准情景围绕East与稳定核心产品配置资源，悲观情景收紧可撤销预算并守住West，旺季压力情景则以Phones上侧风险边界检验销售额承接能力；Central等敏感项只有在评分与实际增量同时满足条件时才追加投入。该分档只调整资源承诺强度，不虚构利润最优、价格弹性或库存件数。
